\chapter{統計的裁定と平均回帰}
\label{ch:statarb}

\section{戦略の中身}

統計的裁定は、似た値動きをする銘柄の間の一時的な「ずれ」が戻ることに賭ける。代表的な形は三つある。

\begin{enumerate}
\item \textbf{ペアトレード}：過去の価格が一緒に動いてきた 2 銘柄の価格差（あるいは比）が、普段より大きく開いたら、上がった方を売り、下がった方を買う。差が戻ったら手じまう。
\item \textbf{ファクターの残差の平均回帰}：各銘柄のリターンを、市場・業種のファクター（主成分や ETF）に回帰し、その残差の累積が平均から離れたら、戻る方向に賭ける。数百〜数千銘柄に同時に小さく賭ける。
\item \textbf{短期の反転}：直近の 1 週間〜1 か月に大きく下げた株を買い、大きく上げた株を売る。
\end{enumerate}

保有期間は数分〜数週間で、回転率が高い。収益の源は、一時的な需給の偏り（大口の売買、指数の資金の流れ、損切り）に対して流動性を供給する対価と解釈されることが多い。

\section{短期の反転}

Lehmann（1990）\cite{lehmann1990} は週次、Jegadeesh（1990）\cite{jegadeesh1990} は月次で、米国株のリターンに強い負の自己相関（反転）があることを示した。直近の負け組を買い勝ち組を売る逆張りの戦略は、取引コストの前では大きな収益を上げた。

\begin{paperbox}[Nagel（2012）「消える流動性」\cite{nagel2012}]
短期の反転の戦略の収益は、流動性を供給する対価と解釈できる。この収益は、市場の変動率（VIX）が高い時期に大きくなる。流動性の供給者（マーケットメイカーやヘッジファンド）自身が資金の制約を受け、流動性の供給が細る時期に、その対価が上がるからである。2008 年の金融危機では、反転の戦略の期待収益が大きく上がった。
\end{paperbox}

ただし、反転の戦略は回転率が非常に高く、取引コストの後にはほとんどの収益が消える（Novy-Marx・Velikov 2016 \cite{novymarxvelikov2016}）。実際にこの収益を取れるのは、コストの低い高頻度の業者と、執行の上手な大手の運用会社に限られる。

\section{ペアトレード}

\begin{paperbox}[Gatev・Goetzmann・Rouwenhorst（2006）\cite{ggr2006}]
米国株（1962〜2002 年）で、12 か月の価格の動きが最も近い（正規化した価格の距離が最小の）組をペアとし、その後 6 か月の間、価格差が過去の標準偏差の 2 倍を超えて開いたら売買する単純な規則を検証した。自己資金の要らないペアのポートフォリオは、年率で最大 11\% 程度の超過リターンを上げ、保守的な取引コストを引いても残った。ただし、収益は期間の後半に小さくなっていた。
\end{paperbox}

Do・Faff（2010）\cite{dofaff2010} は、同じ方法を 2009 年まで延ばし、ペアトレードの収益が 1990 年代以降に低下していること、ただし市場の混乱の時期（2000〜2002 年、2007〜2009 年）には収益が上がることを示した。

ペアを選ぶ方法として、価格の距離のほかに、\textbf{共和分}（Engle・Granger 1987 \cite{englegranger1987}）がよく使われる。二つの価格の系列がそれぞれランダムウォークでも、ある線形結合 $p^A_t - \beta p^B_t$ が定常（平均に戻る）なら、その結合の値を売買の信号にできる。

\section{ファクターの残差の平均回帰}

\begin{paperbox}[Avellaneda・Lee（2010）\cite{avellanedalee2010}]
米国株の各銘柄のリターンを、主成分（相関行列の上位の固有ベクトル）または業種の ETF のリターンに回帰し、残差の累積を平均回帰の過程（Ornstein--Uhlenbeck 過程）としてモデル化した。残差が平均から標準偏差の何倍か離れたら売買する戦略は、1997〜2007 年の検証で高いシャープレシオを示したが、2003 年以降は成績が下がり、2007 年 8 月に大きく損をした。売買高の情報を加えると成績が改善した。
\end{paperbox}

この形は、クオンツのマーケット・ニュートラルの株式の運用の中心に近い。銘柄の数が多いので、基本法則（\ref{sec:flam} 節）の幅が大きく、1 回の賭けの IC が小さくても、全体の情報レシオは高くなりうる。

\section{夜と昼}

Lou・Polk・Skouras（2019）\cite{lps2019} は、米国株のリターンを、夜間（前日の終値から当日の始値）と日中（始値から終値）に分けると、多くのアノマリーが一方に偏ることを示した。モメンタムの収益はほぼすべて夜間に生じ、ほかの多くのアノマリーの収益は日中に生じ、夜間にはむしろ逆に動く。夜間と日中では主な売買の主体（個人と機関投資家）が違い、その「綱引き」が収益の時間の分布を決めている、と解釈している。

\section{2007 年 8 月の quant quake}

\begin{paperbox}[Khandani・Lo（2011）\cite{khandanilo2011}]
2007 年 8 月 6〜9 日、株式のマーケット・ニュートラルのクオンツのファンドの多くが、市場全体が落ち着いている中で、数日のうちに大きな損失を出し、10 日には大きく戻した。著者らは、単純な短期の反転の戦略（Lehmann 型の逆張り）と、ファクターのポートフォリオのリターンで、この出来事を再現した。
\begin{itemize}
\item ある大きなファンドが、別の資産（信用関連など）の損失や資金の必要から、株式のマーケット・ニュートラルのポジションを急いで解消した可能性が高い。
\item 多くのクオンツのファンドが、似たファクター・似たシグナルで、似たポジションを持っていた。一つの解消が、ほかのファンドの損失を生み、それがさらにリスクの上限による解消を呼んだ（解消の連鎖）。
\item 戦略の収益は 1998 年ごろから低下していて、同じ収益を得るために使うレバレッジが大きくなっていた。
\end{itemize}
\end{paperbox}

この出来事の教訓は、「市場全体のリスクがゼロでも、同じ取引をしている人のリスクは消えない」ことである。リスクモデル（\chapref{ch:risk}）は、価格の過去の共分散からリスクを測るが、ほかの運用者のポジションの重なりは、そこに現れない。Goldman Sachs の幹部が当時「25 標準偏差の動きが何日も続いた」と述べたことは、正規分布を前提にしたリスクの物差しが、混雑の解消の前では役に立たないことを示す逸話としてよく引かれる（Deep Research の答え 1 は、この発言を Khandani・Lo の結果として書いていたが、論文の結論ではない）。

\section{個人や研究で再現できるか}

\begin{itemize}
\item 日次の終値によるペアトレードや反転の戦略は、無料の株価のデータで検証できる。ただし、取引コスト（スプレッドとインパクト）と空売りの費用を正しく入れると、ほとんどの収益が消える。
\item 実際の統計的裁定の優位は、日中の細かいデータ、低い取引コスト、速い執行、多数の銘柄の同時の管理にある。個人が実運用で競うのは難しい。
\item 研究として意味があるのは、「どういう時期に収益が上がるか」（変動率が高く流動性が細る時期、Nagel 2012）や、「誰の売買が歪みを生んでいるか」を調べることである。
\end{itemize}
